% 2009 年江苏卷第 22 题：抛物线 y²=2x 与点 A(2,2)、直线 OA。
% 来源 PDF：普通高考/2009/2009江苏.pdf
% SHA-256：64e417d45911d88d77c047a419dcc259e7c7dc779c79d20df2092755c96a4501
% 原页：2009江苏.pdf 第 2 页（A3 三栏）第 3 栏，书页码 532。
% 绘制依据：按原页画坐标轴、抛物线两支、点 A 及线段 OA，并在两轴标出单位刻度 1。
\begin{tikzpicture}[x=0.95cm, y=0.95cm, semithick, every node/.style={inner sep=1.6pt}]
  \draw[->,>=stealth] (-0.7,0) -- (5.6,0) node[below] {$x$};
  \draw[->,>=stealth] (0,-3.4) -- (0,3.5) node[right] {$y$};

  \draw[domain=-3.15:3.15, samples=140, smooth, variable=\t]
    plot ({\t*\t/2}, \t);

  \draw (0,0) -- (2,2);
  \fill (2,2) circle (2.2pt);
  \node[above right] at (2,2) {$A$};

  \draw (1,0) ++(0,-0.09) -- ++(0,0.18);
  \draw (0,1) ++(-0.09,0) -- ++(0.18,0);
  \node[below] at (1,0) {$1$};
  \node[left] at (0,1) {$1$};
  \node[below left] at (0,0) {$O$};
\end{tikzpicture}
